\section{Introducción}

La inducción electromagnética constituye uno de los principios fundamentales del electromagnetismo y de la ingeniería eléctrica moderna. Este fenómeno, descubierto por Michael Faraday en 1831, establece que una variación temporal del flujo magnético a través de un circuito conductor induce una fuerza electromotriz en dicho circuito. La formulación de este principio permitió el desarrollo de numerosos dispositivos eléctricos, entre ellos el transformador, ampliamente utilizado en los sistemas de generación, transmisión y distribución de energía eléctrica.

Un transformador es un dispositivo electromagnético estático que transfiere energía eléctrica entre dos circuitos mediante un campo magnético variable. Su estructura básica está compuesta por un devanado primario y un devanado secundario enrollados sobre un núcleo ferromagnético. Cuando una corriente alterna circula por el devanado primario, se genera un flujo magnético variable en el núcleo que induce una fuerza electromotriz en el devanado secundario, permitiendo elevar o reducir el nivel de tensión de acuerdo con la relación entre el número de espiras de ambos devanados.

El estudio experimental de los transformadores resulta de gran importancia debido a su papel fundamental en los sistemas eléctricos de potencia. La caracterización de parámetros como la relación de transformación, las pérdidas energéticas y la eficiencia permite comprender el comportamiento de estos dispositivos bajo condiciones reales de operación y comparar los resultados experimentales con las predicciones del modelo ideal.
En la presente práctica se analiza el funcionamiento de un transformador  a partir de la Ley de Inducción de Faraday. Para ello, se determina experimentalmente la relación entre los voltajes del devanado primario y secundario, se evalúa la influencia del número de espiras en la transformación de tensión y se calcula la eficiencia energética del sistema bajo diferentes condiciones de carga. Adicionalmente, se estudia el papel del núcleo ferromagnético en la conducción del flujo magnético y su influencia en el desempeño del transformador.
